\section{Where Have All the Heroes Gone}

\begin{quotex}
Unless one is pious, one cannot be truly wise.
\flright{\textsc{Giambattista Vico}}
\end{quotex}


\textbf{Giambattista Vico} in his New Science tried to develop a science of civil society. He pointed out that, unlike nature which was created by God, civil society was the creation of man. Hence, while nature is opaque, civil society should be clear. This may be partly true, but Vico does not take into account the subterranean and higher forces that are also involved in the creation of civil society. Furthermore, even the part played by man is opaque, because men, with rare exceptions, are decidedly lacking in effective self-knowledge. The net result is the opposite of what he expected. Great progress has been made in the science of nature while the social sciences are, for the most part, risible. Anthropologists, for example, have recently given up the pretense to be considered a science.\footnote{\url{http://archaeology.about.com/b/2010/12/15/is-anthropology-a-science-aaafail.htm}}

Nevertheless, Vico's study of history is groundbreaking since he included myths and symbols as part of his raw material. He introduced a developmental scheme by which societies progress from a state of nature to higher levels. Unlike Rousseau's social contract, Vico claims that society in the state of nature was ruled by the fathers. The fathers provided peace, prosperity, protection, and order. Initially, society consisted of the families and their servants or clients. \emph{Strangers, i.e., those without gods, without caste, without even fathers, were attracted to the order offered by the fathers; thus they migrated to those realms, where they remained a constant threat to the peace.}

Vico sees three stages which he attributes to differences in human nature. As we shall see, the differences are actually related to the various castes, and where Vico sees a progress, the men of tradition see a decline or a degeneration of castes. These stages are characterized by \emph{the theological poets, the heroes, and finally reasonable men}.


\paragraph{Theological Poets}


The first stage was dominated by the theological poets, according to Vico. These men had a poetic or creative nature dominated by the imagination, which Vico regarded as delusional. He anticipated the later demythologizers, since he claimed the enlightenment philosophers obtained the same wisdom without the illusions. At this stage, the urges of men were kept in check by their religion. Men believed the source of law was divine and that life and human affairs depended on the gods. The sages were ``initiates'', i.e., the interpreters of the gods. Authority comes from the divine, and clearly, the gods, or God, could not be called into account.

What follows is my interpretation of Vico's findings. These first fathers, then, were priests. As priests, they would interpret God's will, understand the divine law, and offer sacrifices. In spiritual matters, they were infallible, as no one could doubt God's law or His will. Rather than a figment of their imagination, these fathers had a direct intuition of the higher things.

\paragraph{The Heros}


After the priests, the heroes gained control. They believed they descended from the gods and had a natural nobility. Power is what counted for them, as the divine law was replaced by the law of force. When they paid deference to the priest, force was kept within the bounds of duty. Otherwise, rule was by the most powerful, not the most wise.

Clearly, the heroes correspond to the warrior caste. Their knowledge of the transcendent was restricted to the practical, rather than the pure or speculative, realm. Thus, in the world of action they inhabited, they pursued the transcendent virtues of nobility, strength, courage, loyalty, and so on. Unable to intuit God's will directly, they were often subject to the boons and the banes of the gods.

\paragraph{Civilized Men}


Interestingly, Vico calls the next state that of the truly human or civilized men. This stage is intelligent, moderate, benign, and reasonable. Men at this stage are guided by the laws of conscience, reason, and duty. In such societies, there is no longer a caste system, since all men are born free and equal. Law is created by man for his benefit and is no longer God's law nor the edicts of the powerful.

We see this as the rule of the caste of the bourgeoisie. As they are involved in the world of forms, their use and transformation, any appeal to a higher authority is interpreted as an unwelcome restriction on their activity.

\paragraph{Proofs and Refutations}


If Vico believes he discovered an ineluctable law of human development, and not just some contingent events exclusive to Europe, then immediate counterexamples come to mind. First of all, there should be demonstrable examples of undeveloped societies that have not reached the third stage. We do see them, and they are brutal and unjust as he describes. However, what we do not observe are the great esoteric poets and heroes who should be dominating such groups in their undeveloped stage. Rather than undeveloped, they are more likely at a fourth stage, not known to Vico, which we shall discuss presently.

From the traditional point of view, we really do not see any progress; instead, we witness a series of privations. At the first stage of the priest/poets, there was a direct connection with the divine, which was then lost. The next privation was the hierarchy of the heroes with their virtues, which was replaced by human reason, egalitarianism, and the pursuit of comforts. To be fair to Vico, his views were nuanced by the monarchies and the ``true religion'', but that will be for another time.

\paragraph{The Revolt of the Masses}


There was, and is, a further stage which Vico had not anticipated, which is the revolt of the masses, originally the strangers, against the fathers. At this point, the civil society is without priest/poets, without heroes, and now even without reasonable men. Reasonable men are guided truth and justice, but the revolt denies both of these. Instead of ``truth'', thoughts are rated on whether or not they are offensive. Since the mass is necessarily disordered, the variety of offensives is quite large. Rational justice is based on what is rightly due to someone; the justice of the masses is based solely on need, regardless of effort and character.

Contrary to the fantasies of Freudianism, good sons don't revolt against their fathers. They may chafe at such rule in their youth, but later in life they come to praise their fathers. So, in his \emph{Moses and Monotheism}, Freud is not describing the Oedipal patricide, but rather the revolt of the masses against the source of order. This was resisted by the Heroes, but with the influx of the idea of equality, and a distaste for violence, such revolts became possible. Unfortunately for the ``reasonable men'', the masses don't share the same distaste for violence, and they have often become the victims as in France, Russia, China, Cambodia, and so on.

\paragraph{Signs of Order}


Vico points out certain universal customs in the Poetic and Heroic ages that indicate the divine order, to wit:

\begin{itemize}


\item All have some religion

\item All contract solemn marriages

\item All bury their dead
\end{itemize}


Hence, the revolt always and everywhere attacks these customs. Religion is mocked and denied. Next, the marriage bond is weakened and even redefined. Finally, respect and reverence for the dead fades. The masses have no connection, neither a spiritual nor even, in many cases, a genetic, connection to the fathers. Hence, their characters are besmirched.

Do I really need to provide specific examples? Of course, all this is regarded as progress, even as ``scientific''. Didn't Galileo discover the moons of Jupiter? But what analogous scientific discovery will justify such progress?

The assumption is that order will remain even if the sources of order are destroyed. The assumption is that when everyone is tolerant of each other, then unity will arise. The assumption is that, when people are provided with the material goods they lack, then they will be happy. The assumption is that when people are encouraged to satisfy their lowest desires, then they will be satisfied. The assumption is that mind-altering drugs are good, since there is nothing worth thinking about anyway. The assumption is that when the past is forgotten, then the future will be boundless.


\flrightit{Posted on 2013-03-11 by Cologero}

\begin{center}* * *\end{center}

\begin{footnotesize}
\begin{sffamily}
\texttt{Saladin on 2013-03-14 at 20:23 said:}

``Rather than undeveloped, they are more likely at a fourth stage, not known to Vico, which we shall discuss presently.''

We are heading there, if we are not already there!

\hfill

\texttt{Logres on 2013-03-15 at 15:21 said:}

``Even to dispraise Sauron, then, was held to be treason''... .

Tolkien, writing of Numenor's fall

\hfill

\texttt{Izak on 2013-03-17 at 17:27 said:}

Nice to see some discussion on Vico!

I asked for and received the Penguin edition of his New Science for Xmas, so once I get some time to read it, I'll have to revisit this post with a fresher and more solid perspective on what Vico was trying to do.

\hfill

\texttt{Cologero on 2013-03-17 at 21:02 said:}

Those reading the mailing list will understand why the topic of Vico just came up. Vico was probably the first to take ancient symbols and myths seriously as tools for prehistoric research. However, he takes a somewhat evolutionistic view, seeing human history as a progress from inferior states to superior. Nevertheless, there is much good material in his New Science.

\hfill

\texttt{Matt Edge on 2016-03-07 at 06:55 said:}

I love reading this blog, following your thoughts is as exciting as wrestling and I feel them adding strength to my mind. Thank you for this blog.

\hfill

\end{sffamily}
\end{footnotesize}

